% !TeX root = ../main.tex
% !TeX encoding = UTF-8

\chapter{The Standard That Was Already There, Medical Law and the Evidence the Profession Built in 1957}

In October 1954, John Bolam was admitted to Friern Hospital in North London for treatment of a depressive illness. His consultant recommended electroconvulsive therapy. Mr Bolam consented to the treatment. He was not warned that ECT carried a risk of bone fracture, a risk that, at the time, was accepted by a responsible body of medical opinion as an inherent consequence of the procedure that did not require disclosure. During the treatment, Mr Bolam sustained fractures to his pelvis. He sued the hospital authority for negligence.

Justice McNair's judgment in Bolam v Friern Hospital Management Committee {[}1957{]} 1 WLR 582 established the foundational standard for medical negligence liability in English law. A doctor is not negligent if they act in accordance with a practice accepted as proper by a responsible body of medical opinion skilled in the particular form of treatment in question. The operative words, for the purposes of this book, are not ``responsible body of medical opinion.'' They are ``at the time of the decision.''

Bolam established that a doctor's conduct is assessed against what was known and what was accepted practice at the moment the clinical decision was made, not at the moment of trial, not at the moment the harm became apparent, not at the moment a better standard became available. The information picture at the decision moment is the evidentiary universe within which the clinician's conduct is judged. What came after is irrelevant to the question of whether the duty of care was discharged. What existed at the moment is everything.

Medical law arrived at the temporal evidentiary standard this book calls the Digital Alibi sixty-eight years before artificial intelligence made it a governance question for corporate boards. It arrived at it through litigation, refined it through decades of case law, and has now embedded it so deeply in clinical practice that its absence would be professionally and legally inconceivable. The contemporaneous clinical record, the note written at the time of the clinical decision, recording what the clinician knew, what they considered, and what they decided, is not a bureaucratic habit. It is the clinician's alibi.

Chapter 10 is the final case study in Part Two, and it is different from the five that preceded it. Chapters 5 through 9 examined industries in which the Digital Alibi gap produced catastrophe, the Grand Canyon collision, the 2008 financial crisis, the Post Office prosecutions, the Challenger disaster, the Lloyd's Names ruin. This chapter examines an industry that built the Digital Alibi standard, progressively and through hard experience, before the catastrophe that would have made it unavoidable. Not fully. Not without remaining gaps. But sufficiently to demonstrate what a profession looks like when the evidentiary standard is not aspirational but operational.

\scenebreak

\section{Thalidomide and the Regulatory Reckoning: The Information Picture That Was Not Required}

\textbf{The decision without the evidence} Frances Kelsey, FDA, 1960-1962

Thalidomide was approved as a sedative and anti-nausea drug in West Germany in 1957. It was subsequently licensed in forty-six countries and marketed widely for use in pregnancy. Between 1957 and 1962, an estimated ten thousand children were born with severe limb deformities, phocomelia, as a direct consequence of their mothers' thalidomide use during the first trimester of pregnancy.

The thalidomide disaster is, for the purposes of this book, a case about the information picture that regulatory approval did not require. The licensing decisions that allowed thalidomide into the market were made on clinical trial data that did not include adequate teratogenicity testing, testing of the drug's effects on foetal development. That testing had not been required. The information picture at each national approval decision did not include what the drug did to developing foetuses because nobody had yet established that it was necessary to find out.

Frances Kelsey, the FDA medical officer responsible for reviewing Richardson-Merrell's application for US market approval, repeatedly refused to approve thalidomide on the grounds that the clinical trial data submitted was insufficient to establish safety. She required what was not in the application: evidence adequate to answer the questions that needed answering, not evidence adequate to satisfy the requirements as they then existed. The information picture at the decision moment, she concluded, did not justify approval. The US did not license thalidomide. The limb deformities that affected European and Australian children did not occur at scale in the United States.

The regulatory response to thalidomide, in the United States through the 1962 Kefauver Harris Amendment, in the United Kingdom through the Committee on Safety of Drugs established in 1963, created the clinical trial documentation requirements that now govern pharmaceutical development worldwide. Those requirements do not merely specify what tests must be conducted. They specify what information must be captured, in what form, at what stage of the trial, in a contemporaneous record that is available for regulatory review. The information picture at each stage of the clinical decision process, from preclinical testing through phase three trials to post-marketing surveillance, must be documented as it exists, not assembled retrospectively when the application is reviewed.

\begin{quote}
\emph{Frances Kelsey did not refuse thalidomide approval because the drug was unsafe. She refused it because the information picture at the approval decision moment did not contain adequate evidence to establish that it was safe. The distinction between ``we have not found it unsafe'' and ``we have evidence that it is safe'' is the evidentiary standard she imposed. It is the Digital Alibi standard.}
\end{quote}

\scenebreak

\section{From Bolam to Montgomery: Sixty Years of Tightening the Standard}

The Bolam test, as established in 1957, assessed clinical conduct against the standard of a responsible body of medical opinion. A doctor who could demonstrate that their practice was consistent with what responsible colleagues would have done had established, in principle, that the duty of care had been discharged. The information picture at the decision moment was assessed by reference to what responsible medical opinion would have known and done, not by reference to a fixed external standard.

For nearly sixty years, Bolam governed medical negligence in English law. During that period, the clinical record, the contemporaneous note of what the clinician knew and decided, was the primary evidentiary instrument. A clinician who had made a contemporaneous record of their clinical reasoning, the information available to them, and the decision they reached was, in most cases, in a defensible evidential position. A clinician who had not made that record was not.

The Montgomery v Lanarkshire Health Board {[}2015{]} UKSC 11 judgment of the UK Supreme Court fundamentally tightened the standard for informed consent. The Bolam test had allowed a clinician to discharge the duty to disclose risks by demonstrating that a responsible body of medical opinion would not have disclosed them. Montgomery replaced this with a patient-centred standard: a clinician must take reasonable care to ensure that the patient is aware of any material risks involved in the proposed treatment, and of any reasonable alternatives. A material risk is one to which a reasonable person in the patient's position would be likely to attach significance. The clinician's assessment of what a responsible body of colleagues would have disclosed is no longer the standard. The patient's informational position at the moment of consent is the standard.

Montgomery's significance for the Digital Alibi argument is precise. Bolam asked: did the clinician's conduct conform to responsible practice? Montgomery asks: was the patient's information position at the consent moment adequate for genuine informed decision-making? The shift from the clinician's practice standard to the patient's information position is the shift from a process question to an evidentiary question. It is the same shift this book makes for AI governance: not ``did the board have a governance process?'' but ``was the board's information position at the decision moment adequate to discharge its duty of care?''

\begin{quote}
\emph{Bolam asked whether the clinician's conduct was proper. Montgomery asked whether the patient's information position at the consent moment was adequate. That shift, from process standard to evidentiary standard, is the same shift this book makes for AI governance. The question is not whether the board had governance. The question is whether the information picture at the decision moment was adequate to discharge the duty of care.}
\end{quote}

The practical consequence of Montgomery for clinical practice was the intensification of contemporaneous consent documentation. A clinician who cannot produce a contemporaneous record of what information was given to the patient, what the patient's questions were, and what the patient indicated they understood before consenting is not in a defensible evidential position after Montgomery. The contemporaneous consent record is not optional documentation. It is the evidence that the patient's information position at the consent moment met the legal standard.

\scenebreak

\section{The Contemporaneous Clinical Record: The Digital Alibi in Practice}

The contemporaneous clinical record that medical law has made mandatory is, in its evidentiary architecture, a complete implementation of the Digital Alibi protocol for clinical decisions. Its four elements map precisely onto the four Digital Alibi elements this book has defined.

Decision Identity. The clinical record identifies the clinician making each decision, their professional qualifications and authority, the date and time of the decision, and the specific clinical action taken. A clinical note that records ``Dr.~Smith, Consultant Cardiologist, reviewed patient at 14:32 on 3 March 2024 and decided to adjust medication regime'' has established Decision Identity with the precision the Digital Alibi requires. It names who made the decision, in what capacity, with what authority, at what exact moment.

The Information Picture. The clinical record documents the clinical information available at the decision moment: the patient's presenting symptoms, the relevant history, the test results available, the differential diagnoses considered, and, crucially, after Montgomery, the information given to the patient about risks and alternatives. The negative space is also required: a note recording that a specific investigation had not yet returned, or that the patient had not disclosed a relevant condition, captures what was not known at the decision moment. The clinical record is required to reflect the information picture as it existed, not as it was subsequently found to be.

The Human Oversight Record. The clinical record documents the oversight exercised at the decision moment: who reviewed the patient, in what capacity, with what clinical authority. Where a decision involves multiple clinicians, a team discussion, a multidisciplinary meeting, a second opinion, the record documents who was present, what was discussed, and what each participant contributed to the decision. The oversight record is not a summary of the team's eventual conclusion. It is a record of who did what, when, and on what basis.

The Accountability Chain. The clinical record establishes named individual accountability from bedside to board in a form that has been legally tested and refined over decades. The treating clinician is named. The supervising consultant is named. The clinical governance structure within which the decision was made is documented. Where accountability transfers, between shifts, between specialties, between inpatient and outpatient settings, the transfer is recorded. The chain is contemporaneous, named, and specific.

\begin{longtable}[]{@{}
  >{\raggedright\arraybackslash}p{(\columnwidth - 0\tabcolsep) * \real{1.0000}}@{}}
\toprule\noalign{}
\begin{minipage}[b]{\linewidth}\raggedright
\textbf{THE CLINICAL RECORD AS DIGITAL ALIBI, A COMPLETE IMPLEMENTATION}
\end{minipage} \\
\midrule\noalign{}
\endhead
\bottomrule\noalign{}
\endlastfoot
Decision Identity: clinician named, qualified, and authorised; decision timestamped; specific action recorded. Not attributed to a team or a process, attributed to a named individual at a named moment. \\
Information Picture: presenting information, available test results, history, differential diagnoses, and the negative space (what had not yet returned, what the patient had not disclosed) all documented at the decision moment. \\
Human Oversight Record: reviewers named, their capacity and authority recorded, the substance of their review documented, not that a review occurred, but what it consisted of and what it concluded. \\
Accountability Chain: named individual accountability from treating clinician to supervising consultant to clinical governance structure, documented contemporaneously, with transfers of accountability recorded at each handoff. \\
Montgomery compliance: contemporaneous consent record documenting what information was given, what the patient's questions were, and what the patient indicated they understood, not because it is good practice, but because it is the evidence that the patient's information position met the legal standard. \\
The Digital Alibi lesson: this architecture is not unique to medicine. It was built in response to litigation that established the same evidentiary standard this book applies to AI-assisted decisions. Every industry arrives at this standard eventually. Medicine arrived at it first. \\
\end{longtable}

\scenebreak

\section{What Medicine Has Not Yet Solved, and Why It Matters for AI}

Clinical medicine has built the most developed Digital Alibi architecture of any profession. The contemporaneous clinical record, refined over decades of litigation and regulatory development, captures the four elements with greater systematic rigour than any other governance framework in any sector. But it has not fully solved the problem. Two remaining gaps in clinical governance are directly relevant to AI governance, because AI-assisted clinical decisions are creating versions of both.

The first gap is the negative space at system level. The clinical record captures the negative space for individual clinical encounters, what was not known, what had not returned, what the patient had not disclosed. It does not, in most clinical environments, capture the systemic negative space: what patient populations the clinical guidelines were not designed for, what research the evidence base had not yet incorporated, what structural biases existed in the clinical training data that informed the guideline development. An AI clinical decision support system that operates on guidelines built from a research base that excluded certain patient populations has a systemic negative space that neither the clinical record nor the current regulatory framework is designed to capture. The Montgomery standard requires what was not disclosed at the individual consent moment to be documented. It does not require what the clinical system could not see to be documented at the population level.

\begin{longtable}[]{@{}
  >{\raggedright\arraybackslash}p{(\columnwidth - 0\tabcolsep) * \real{1.0000}}@{}}
\toprule\noalign{}
\begin{minipage}[b]{\linewidth}\raggedright
\textbf{THE BOARDROOM MIRROR, AI DIAGNOSTIC TRIAGE AND THE MEDICAL GOVERNANCE GAP}
\end{minipage} \\
\midrule\noalign{}
\endhead
\bottomrule\noalign{}
\endlastfoot
The scenario: an NHS trust deploys an AI triage system to assist emergency department clinicians in prioritising patients by clinical urgency. The system is validated on retrospective patient data, approved through the trust's clinical governance committee, and deployed with a policy position that it supports, not replaces, clinical judgement. The clinical record for each patient documents the clinician's assessment and decision. It documents that the AI system's output was considered. \\
The systemic negative space: the AI triage system was trained on historical emergency department attendance data from a trust whose catchment area is predominantly urban. Its training data contains a structural underrepresentation of patients presenting with atypical symptom profiles associated with specific demographic groups, older women with cardiac events, patients with learning disabilities presenting pain indirectly, non-English speakers whose symptom descriptions were less precisely captured in historical records. The system's triage outputs systematically underweight urgency for these presentations. No individual clinical record captures this pattern. Each record documents a clinical assessment. None documents what the AI system could not see about this patient's presentation type. \\
The information picture gap: the clinical record shows the model was consulted and clinical judgement was applied. It does not show, for each specific patient assessment, what the AI system was not given, which features of that patient's presentation were structurally underweighted by the training data distribution. The negative space exists at the population level, invisible in any individual record. \\
The Montgomery parallel: Montgomery requires that the patient's information position at the consent moment be adequate for genuine informed decision-making. A patient triaged by an AI system with a structural blind spot to their presentation type has not been triaged on an adequate information picture, but no clinical record documents that inadequacy, because the clinical record was designed for human decision-making and captures what the clinician knew, not what the AI system could not see. \\
The Digital Alibi lesson: even the most developed governance standard in any profession does not yet fully capture the AI-specific dimension of the information picture gap. The contemporaneous clinical record is the gold standard for human clinical decisions. For AI-assisted decisions, it requires extension: the output-moment capture of what the system was not given, at the level of each specific assessment, must be added to the record architecture that medicine built for human decisions alone. \\
\end{longtable}

The second gap is the algorithmic accountability chain. Where clinical decisions are supported by AI systems, diagnostic imaging analysis, clinical risk stratification, treatment recommendation engines, the accountability chain from algorithm to clinician to patient is not yet consistently documented in the clinical record with the precision that accountability chains for purely human decisions require. Who is accountable for the version of the algorithm deployed at the moment of each specific clinical decision? What information did the algorithm process? What was it not given? These questions are not yet routinely answered by the clinical record for AI-assisted decisions, because the clinical record's architecture was designed for human decision-making and has not been fully extended to capture the additional elements that AI-assisted decision-making requires.

\begin{quote}
\emph{Medicine has built the most complete Digital Alibi architecture of any profession. Its remaining gaps are precisely the gaps that AI creates: the systemic negative space the record was not designed to capture, and the algorithmic accountability chain the record's architecture does not yet reach.}
\end{quote}

\scenebreak

\section{What Part Two Has Established}

Six industries. Six distinct forms of the same failure. Six governance architectures built in response, each of which is, in its essential design, a Digital Alibi protocol.

Aviation built the Flight Data Recorder mandate because the Grand Canyon collision demonstrated that the information picture at the moment of a consequential aviation decision must be forensically reconstructible, permanently, in real time, or the causes of catastrophic failure cannot be established and the lessons cannot be learned. The three design principles it embedded, continuity, independence, and survivability, are the design principles of any Digital Alibi architecture.

Financial services built the SM\&CR and model risk management framework because the 2008 crisis demonstrated that institutional accountability without named individual accountability at the decision moment allowed catastrophic governance failures to produce institutional consequences without individual ones. The shadow accountability gap, the gap between the formal record and the operative decision channel, is the same gap that AI governance's approval culture creates.

The Post Office Horizon scandal demonstrated, in the most consequential form the Digital Alibi gap can take, what happens when a policy position is maintained without a contemporaneous information picture at each decision moment: an injustice can be sustained for decades by institutional momentum, because the contemporaneous record that would have forced the reckoning was never created.

The Challenger disaster demonstrated that the Digital Alibi gap is not always about absence. It can be about distortion, a governance process that structures the information picture in a way that makes a correct recommendation appear unsupportable. The burden-of-proof inversion that produced Challenger is replicated by AI approval culture whenever a governance architecture assumes outputs are correct unless proven otherwise.

Lloyd's of London demonstrated that the Digital Alibi gap is not only a real-time problem. Long-tail liabilities, asbestosis, pollution, reinsurance spirals, arrive decades after the decision moments that created them. The information picture at each of those decision moments must survive not just immediate scrutiny but the full duration of the exposure. The Legacy Poisoning risk in AI governance is the same long-tail threat in a new technological domain.

And medicine demonstrated that the Digital Alibi standard is not a future aspiration. It is an existing requirement, built through litigation, embedded through regulation, and operational in every clinical encounter in England and Wales. Medicine arrived at it through catastrophe, Thalidomide, the cases that refined Bolam, the patient whose information position at the consent moment was found inadequate in Montgomery. Other industries arrived at their equivalents through their own reckonings. The boards reading this book have not yet had their reckoning. Part Three explains, in operational detail, what their reckoning will look like when it arrives.

\begin{quote}
\emph{Medicine did not build the contemporaneous clinical record because clinicians are more virtuous than other professionals. It built it because litigation established, case by case, that the information picture at the decision moment is the evidentiary universe within which conduct is judged. Every industry arrives at this standard. The only question is whether they arrive at it before or after the cases that establish it.}
\end{quote}

\scenebreak

\section{From History to Architecture: Part Three}

Part Two has examined six industries. In each of them, the Digital Alibi gap produced a specific category of consequence, mandate, settlement, prosecution, disaster, ruin, and, in medicine, an ongoing refinement of the standard that makes the gap visible as it narrows. In each of them, the governance architecture built in response is a version of the same protocol: capture the information picture at the decision moment, in a form that is contemporaneous, independent, and survivable.

Part Three applies this architecture to the specific operational reality of AI-assisted decisions in regulated organisations today. It takes the four Digital Alibi elements established in Chapter 2, the legal framework established in Chapter 3, and the structural analysis of existing frameworks established in Chapter 4, and asks: what does building the Digital Alibi architecture actually require, system by system, decision by decision, in the organisations where AI is already deployed?

The answer is operational, not theoretical. Part Three is the bridge between the argument this book has made and the protocol Part Four provides. By the time a board finishes Part Three, they will know precisely where their organisation sits in the governance gap, what the four elements require in operational terms for each of their material AI systems, and what the specific liability exposure of their current position is. The argument will have produced its diagnosis. Part Four provides the prescription.
